\FloatBarrier
\section{Data samples}
\label{Sec:Model}
%PYTHIA
%- Overview of PYTHIA
%- How is hadronisation described in PYTHIA
%- Tunes
%    - MONASH
%    - JUNCTIONS
%    (we could also look at:
%    - Close-packing
%    - Ropes)
%Table with number of events, number of generated b and c quarks, yield of B+, B-, $\Lambda_b$ (and bar), D+, D-, $\Lambda_c$ (and bar).
%Multiplicity spectra
%Some kinematic spectra (pT, $\eta$, $\phi$, etc.) for the particles we care about. We can move this to some appendix but it will be nice to have here as an overview.


This section describes the simulated pp samples used in the analysis, the common particle-level selections, and the three PYTHIA~8 hadronisation configurations compared throughout the paper: MONASH, JUNCTIONS, and CLOSEPACKING.

\subsection{PYTHIA~8 setup}
\label{sec:pythia-setup}

Events are generated with PYTHIA~8 using the Lund string model for hadronisation.  Proton--proton collisions are generated at $\sqrt{s} = 14~\mathrm{TeV}$. The samples are hard-heavy-flavour enriched samples rather than minimum-bias samples.  
The common generator settings are summarized in Table~\ref{tab:generator-settings}.

\begin{table}[tbp]
\centering
\label{tab:generator-settings}
\begin{tabular}{ll}
\toprule
Setting & Value \\
\midrule
Collision system & pp \\
Centre-of-mass energy & $\sqrt{s}=14~\mathrm{TeV}$ \\
Beam IDs & \texttt{Beams:idA = 2212}, \texttt{Beams:idB = 2212} \\
Baseline tune & \texttt{Tune:pp = 14} \\
Hard charm production & \texttt{HardQCD:hardccbar = on} \\
Hard beauty production & \texttt{HardQCD:hardbbbar = on} \\
Hard-process cut & \texttt{PhaseSpace:pTHatMin = 1.} \\
Events per job & $10^{6}$ \\
Jobs per tune & $100$ \\
Events per tune & $10^{8}$ \\
\bottomrule
\end{tabular}
\caption{Common PYTHIA~8 generator settings used for all three hadronisation configurations.}
\end{table}

\subsection{Particle selection}
\label{sec:particle-selection}

Particles are selected at generator level after hadronisation.  The analysis uses final-state particles satisfying $p_{\mathrm{T}} \geq 0.15~\mathrm{GeV}/c$ and $|\eta| \leq 4$. The weakly decaying heavy-flavour hadrons used in the analysis are kept stable at generator level by disabling their decays.

The charged-particle multiplicity, $N_{\mathrm{ch}}$, is defined using prompt charged particles in the same kinematic acceptance with PYTHIA status codes in the prompt range used by the analysis.

\subsection{Hadronisation models}
\label{sec:hadronisation-models}

The same hard-process setup and particle-level selections are used for all three configurations.  The configurations differ only in the modelling of hadronisation, colour reconnection, and dense-string effects.

\subsubsection{MONASH}
\label{sec:monash}

The MONASH configuration is used as the baseline PYTHIA~8 reference tune.  It corresponds to the Monash 2013 tune, which provides a global tuning of fragmentation, multiparton interactions, colour reconnection, and underlying-event parameters.

This configuration is used as the reference because it is the standard PYTHIA~8 pp tune and provides the baseline expectation for heavy-flavour hadronisation without the additional junction or close-packing mechanisms studied here.

\subsubsection{JUNCTIONS}
\label{sec:junctions}

The JUNCTIONS configuration uses the QCD-based colour-reconnection model with junction topologies enabled.  In this model, colour reconnections are guided by QCD colour rules and by minimization of the string-length measure. 

\subsubsection{CLOSEPACKING}
\label{sec:closepacking}

The CLOSEPACKING configuration includes the string close-packing model.  The model accounts for the effect of overlapping strings in high-density pp events.  Close strings modify the effective string tension, which can affect strangeness production, transverse-momentum spectra, and baryon production.

\subsection{Generated samples}
\label{sec:generated-samples}

For each configuration, $100$ independent jobs are generated, each containing $10^{6}$ events.  The total generated sample therefore contains $10^{8}$ events per configuration and $3\times10^{8}$ events overall.  The generated sample sizes and selected heavy-flavour hadron yields are shown in Table~\ref{tab:generated-heavy-flavour-summary}.

\begin{table}[tbp]
\centering
\label{tab:generated-heavy-flavour-summary}
\begin{tabular}{lrrrr}
\toprule
Observable & MONASH & JUNCTIONS & CLOSEPACKING & Overall \\
\midrule
$N_{\mathrm{ev}}$ & $100\,000\,000$ & $100\,000\,000$ & $100\,000\,000$ & $300\,000\,000$ \\
$N_{b+\bar b}^{\mathrm{acc}}$ & $13\,321\,922$ & $13\,258\,188$ & $13\,263\,640$ & $39\,843\,750$ \\
$N_{c+\bar c}^{\mathrm{acc}}$ & $125\,387\,414$ & $124\,842\,736$ & $125\,437\,031$ & $375\,667\,181$ \\
\midrule
$B^{+}$ & $2\,853\,111$ & $2\,044\,989$ & $2\,004\,038$ & $6\,902\,138$ \\
$B^{-}$ & $2\,841\,524$ & $2\,036\,333$ & $2\,004\,631$ & $6\,882\,488$ \\
$\Lambda_{b}^{0}$ & $219\,087$ & $472\,196$ & $973\,428$ & $1\,664\,711$ \\
$\overline{\Lambda}_{b}^{0}$ & $208\,457$ & $474\,179$ & $982\,931$ & $1\,665\,567$ \\
$D^{+}$ & $18\,422\,929$ & $16\,436\,770$ & $16\,405\,528$ & $51\,265\,227$ \\
$D^{-}$ & $18\,474\,130$ & $16\,514\,161$ & $16\,493\,255$ & $51\,481\,546$ \\
$\Lambda_{c}^{+}$ & $2\,282\,077$ & $8\,039\,633$ & $7\,296\,136$ & $17\,617\,846$ \\
$\overline{\Lambda}_{c}^{-}$ & $2\,238\,580$ & $7\,887\,369$ & $7\,195\,995$ & $17\,321\,944$ \\
\bottomrule
\end{tabular}
\caption{Generated-event sample sizes and selected heavy-flavour hadron yields for the three PYTHIA~8 configurations.}
\end{table}

\subsection{Multiplicity classes}
\label{sec:multiplicity-classes}

Event activity classes are defined using percentiles of the charged-particle multiplicity distribution, $N_{\mathrm{ch}}$.  The classes used throughout the analysis are
\begin{equation}
0\text{--}1\%,\ 
1\text{--}10\%,\ 
10\text{--}20\%,\ 
20\text{--}30\%,\ 
30\text{--}40\%,\ 
40\text{--}50\%,\ 
50\text{--}60\%,\ 
60\text{--}70\%,\ 
70\text{--}80\%,\ 
80\text{--}90\%,\ 
90\text{--}100\%.
\end{equation}

The $0$--$1\%$ class corresponds to the highest-multiplicity events, while the $90$--$100\%$ class corresponds to the lowest-multiplicity events.

\begin{figure}[tbp]
    \centering
    \includegraphics[width=0.85\linewidth]{figures/Kinematic Plots/MultiplicitySpectrum_Shared_shape.png}
    \caption{Charged-particle multiplicity distribution for the generated hard-heavy-flavour event samples. The charged-particle multiplicity $N_{\mathrm{ch}}$ is defined using prompt charged particles with $p_{\mathrm{T}}\geq0.15~\mathrm{GeV}/c$ and $|\eta|\leq4$ in pp collisions at $\sqrt{s}=14~\mathrm{TeV}$.}
    \label{fig:mult_distribution}
\end{figure}

\FloatBarrier